\section{Theory}

An unloaded resistive divider (\cref{fig:divider}) produces
\begin{equation}\label{eq:divider}
  V_\mathrm{out} = V_\mathrm{in}\,\frac{R_2}{R_1+R_2}.
\end{equation}
With a load $R_L$ in parallel with $R_2$, replace $R_2$ by $R_2\parallel R_L = R_2R_L/(R_2+R_L)$ in \cref{eq:divider}.

\begin{figure}[ht]
  \centering
  \begin{circuitikz}[scale=0.9, transform shape]
    \draw (0,0) to[V, l=$V_\mathrm{in}$] (0,3)
          to[R, l=$R_1$] (3,3)
          to[R, l=$R_2$, -*] (3,0) -- (0,0);
    \draw (3,3) to[short, -o] (4.2,3) node[right]{$V_\mathrm{out}$};
  \end{circuitikz}
  \caption{Resistive voltage divider.}\label{fig:divider}
\end{figure}

\begin{definitionbox}[Loading effect]
A measuring instrument with input resistance $R_L$ draws current from the circuit and lowers the
quantity it measures; the effect is negligible only when $R_L\gg R_2$.
\end{definitionbox}
